\documentclass[conference,a4paper,flushend]{cs-techrep} % (based on IEEEtran.cls)

% The class iaria.cls loads biblatex/biber with correct IARIA settings
% as well as a set of common packages (times, inputenc[utf8], fontenc[T1],
% graphicx, xcolor, url, orcidlink, hyperref, extdash[shortcuts])

% Packages (self-added, not included in cs-techrep.cls):
\usepackage[shortlabels]{enumitem} % Verbesserte Einstellungsmöglichkeiten von Listen. Nicht bei Beamer anwendbar
% \setlist{itemsep=1ex, topsep=1em}
% \usepackage{graphicx} % Required for inserting images
\usepackage{csquotes} % Korrekte Anführungszeichen setzen. Sprachspezifisch möglich. mit \setquotestyle{style} und \enquote{}.
% \usepackage[american]{babel} % Ändert die Darstellung des Datums in amerikanische Darstellung.
\usepackage{hyperref}
\usepackage[capitalise,nameinlink,noabbrev]{cleveref} % Referenzen besser machen. z.B. mit \cref{}. needs hyperref defined before it.
% \usepackage[backend=biber, style=ieee, alldates=iso]{biblatex} % Zitieren
% \usepackage{array} % Tabellen mit \begin{tabular}
\usepackage{multirow} % Mehrere Zeilen in eine zusammenfassen. für Spalten einfach \multicolumn benutzen
\usepackage{booktabs} % Hübschere Tabellen mit \top/mid/cmid/bottomrule
\usepackage{minted} % Neuere Bibliothek für Code-Blöcke, mit \begin{listing oder minted}
\usepackage{amsmath} % American Mathematical Society für mathematische Formeln, z.B. mit \begin{equation}
\usepackage{amssymb}
\usepackage{amsfonts}
\usepackage{listings}
\usepackage{float}
\usepackage{enumitem}
\usepackage{siunitx}
\usepackage{array}
\usepackage{tabularx}

\addbibresource{references.bib}
% \addbibresource{xampl.bib} % (xampl file provided by the BibTEX base files)
% ======================================================================
% EDIT THESE:
\cstechrepAuthorListTex{Simon Völkl, Erik Dombrowsky, Simon Frey\,\orcidlink{0000-0002-1175-2668}}
\cstechrepAuthorListBib{Simon Völkl and Erik Dombrowsky and Simon Frey}
% Capitalization: https://capitalizemytitle.com/style/Chicago/
\cstechrepTitleTex{LumenTrace: A high-performance MCU-based LED racing simulator} % TODO: title revision?
 % IF you need manual linebreaks in the title, then clone the title without linebreaks for BibTeX:
\cstechrepTitleBib{{\cstechrepTitleTex}}
\cstechrepDepartment{Department of Electrical Engineering, Media, and Computer Science}
\cstechrepInstitution{Ostbayerische Technische Hochschule Amberg\-/Weiden}
\cstechrepAddress{Amberg, Germany}
\cstechrepSeries{Technical Reports}
\cstechrepYear{2026}
\cstechrepMonth{7}
\cstechrepNumber{CL-\cstechrepYear{}-42}
\cstechrepLang{english}  % en-US
% ======================================================================
% DO NOT DELETE AND DO NOT MODIFY THIS:
\filecontentsForceExpansion|[] % force command expansion inside a filecontents* environment
\begin{filecontents*}[overwrite]{selfref.bib}
    @TECHREPORT{selfref,
        author = {|cstechrepAuthorListBib},
        title  = {\cstechrepTitleBib},
        institution = {\cstechrepInstitution, \cstechrepDepartment},
        series = {\cstechrepSeries},
        number = {\cstechrepNumber},
        year   = {|cstechrepYear},
        month  = {|cstechrepMonth},
        langid  = {|cstechrepLang},
    }
\end{filecontents*}

\begin{document}
\selectlanguage{\cstechrepLang}
\maketitle
\begin{abstract}
    % TODO: abstract
    Background: tbd
    Methods: tbd

    Conclusion: tbd
\end{abstract}
\begin{IEEEkeywords}
    MCU, LED, neonpixel, WS2812, racing simulator, Raspberry Pi % TODO: keywords
\end{IEEEkeywords}

% STRUCTURE OF THE TECHNICAL REPORT:
% 1. Introduction
% Motivation: Wieso klassische Slot-Cars durch LEDs ersetzen? (Immersion, Flexibilität, Hardwareverschleiß).
% Unique Selling Point / Contributions: Was macht LumenTrace besonders? (High-Speed Serial Comms, Echtes Physics-Model auf LED-Basis).
% Outline: Kurzer Überblick über den Rest des Reports.

% 2. Related Work
% Welche ähnlichen Ansätze gibt es bereits? (z.B. Open-Source LED Strips Games, kommerzielle digitale Slotcar-Systeme wie Scalextric/Carrera Digital).

% 3. System Architecture & Requirements
% (Hierhin wandern die Requirements, aber deutlich technischer formuliert!)
% Funktionale und nicht-funktionale Systemanforderungen (z.B. Echtzeit-Fähigkeit, 60fps Rendering).
% High-Level Architektur (ähnlich deinem architecture.drawio und der Grafik im README): Trennung von Game Logic (RPI), Controller-Eingaben und und LED-Rendering (Pico).

% 4. Hardware Implementation
% Microcontroller-Setup: Raspberry Pi (4) als Host, Raspberry Pi Pico als MCU für Hardware-IO.
% Kommunikation: UART Protocol (Latenzen, Baud-Raten, Datenpakete).
% Signaltechnik & Controller-Aufbau.

% 5. Software Implementation & Game Mechanics (basierend auf README)
% Core Loop: Game-Tick-Implementierung, Speed Update & Limits.
% Physics & Lane Handling: Reibung, Beschleunigung, Lane-Hops und Kollisionserkennung.
% Dev/Simulation Environment: Der Terminal-Emulator in Python (simulation.py), warum er essenziell für die Logik-Validierung war, bevor geflasht wurde.

% 6. Evaluation (Performance & Results)
% Performance-Auswertung: Wie schnell läuft der Game Loop? (z.B. Python Performance auf dem RPI4).
% Latenz: Delay zwischen Knopfdruck am Controller über den Pico zum RPI und zurück zur LED-Aktualisierung.
% Herausforderungen: Reibungsverluste bei schnellen Updates oder Tearing/Timing-Probleme bei WS2812-Ansteuerung.

% 7. Conclusion & Future Work
% Zusammenfassung des Erreichten.
% Was steht noch auf der Todo-Liste? (z.B. mehr Track-Module, WiFi/MQTT Leaderboards).

\section{Introduction}\label{sec:introduction}

\subsection{Motivation}\label{sec:motivation}
% Wieso klassische Slot-Cars durch LEDs ersetzen? (Immersion, Flexibilität, Hardwareverschleiß).
tbd

\subsection{Unique Selling Point \& Contributions}\label{sec:unique_selling_point}
% Was macht LumenTrace besonders? (High-Speed Serial Comms, Echtes Physics-Model auf LED-Basis).
tbd

\subsection{Outline}\label{sec:outline}
% Kurzer Überblick über den Rest des Reports.
tbd

\section{Related Work}\label{sec:related_work}
% Welche ähnlichen Ansätze gibt es bereits? (z.B. Open-Source LED Strips Games, kommerzielle digitale Slotcar-Systeme wie Scalextric/Carrera Digital).
tbd

\section{System Architecture \& Requirements}\label{sec:system_architecture}

\subsection{Functional and Non-Functional Requirements}\label{sec:requirements}
% Funktionale und nicht-funktionale Systemanforderungen (z.B. Echtzeit-Fähigkeit, 60fps Rendering).
tbd

\subsection{High-Level Architecture}\label{sec:architecture}
% Trennung von Game Logic (RPI), Controller-Eingaben und und LED-Rendering (Pico).
tbd

\section{Hardware Implementation}\label{sec:hardware_implementation}

\subsection{Microcontroller Setup}\label{sec:mcu_setup}
% Raspberry Pi (4) als Host, Raspberry Pi Pico als MCU für Hardware-IO.
tbd

\subsection{Communication}\label{sec:communication}
% UART Protocol (Latenzen, Baud-Raten, Datenpakete).
tbd

\subsection{Signal Processing \& Controller}\label{sec:signal_processing}
% Signaltechnik & Controller-Aufbau.
tbd

\section{Software Implementation \& Game Mechanics}\label{sec:software_implementation}

\subsection{Core Loop}\label{sec:core_loop}
% Game-Tick-Implementierung, Speed Update & Limits.
tbd

\subsection{Physics \& Lane Handling}\label{sec:physics}
% Reibung, Beschleunigung, Lane-Hops und Kollisionserkennung.
tbd

\subsection{Dev/Simulation Environment}\label{sec:simulation}
% Der Terminal-Emulator in Python (simulation.py), warum er essenziell für die Logik-Validierung war, bevor geflasht wurde.
tbd

\section{Evaluation}\label{sec:evaluation}

\subsection{Performance}\label{sec:performance}
% Wie schnell läuft der Game Loop? (z.B. Python Performance auf dem RPI4).
tbd

\subsection{Latency}\label{sec:latency}
% Delay zwischen Knopfdruck am Controller über den Pico zum RPI und zurück zur LED-Aktualisierung.
tbd

\subsection{Challenges}\label{sec:challenges}
% Reibungsverluste bei schnellen Updates oder Tearing/Timing-Probleme bei WS2812-Ansteuerung.
tbd

\section{Conclusion and Future Work}\label{sec:conclusion}
% Zusammenfassung des Erreichten.
% Was steht noch auf der Todo-Liste? (z.B. mehr Track-Module, WiFi/MQTT Leaderboards).
tbd

% ======================================================================
% TEMPLATES FÜR LATEX ELEMENTE
% ======================================================================
% --- Description Template ---
% \begin{description}
%     \item[Item 1] Description 1
%     \item[Item 2] Description 2
% \end{description}
%
% --- Itemize/List Template ---
% \begin{itemize}
%     \item Item 1
%     \item Item 2
% \end{itemize}
%
% --- Listing/Code Template ---
% \begin{lstlisting}[language=json, float, caption={example listing}, label={lst:listing}]
% {}
% \end{lstlisting}
%
% --- Figure Template ---
% \begin{figure*}[!htb]
%   \centering
%   \includegraphics[width=\linewidth]{img}
%   \caption{caption}
%   \label{fig:figure}
% \end{figure*}
%
% --- Table Template 1 (Standard with Tabularx) ---
% \begin{table}[!h]
%     \centering
%     \caption{Table Caption}
%     \label{tab:table_label1}
%     \begin{tabularx}{0.9\linewidth}{>{\raggedright\arraybackslash}p{2.8cm} >{\raggedright\arraybackslash}X}
%         \toprule
%         \textbf{Category} & \textbf{Content} \\
%         \midrule
%         yada & yada \\
%         \bottomrule
%     \end{tabularx}
% \end{table}
%
% --- Table Template 2 (Full Width Costs Table) ---
% \begin{table*}[!htb]
%     \centering
%     \caption{Costs}
%     \label{tab:costs}
%     \begin{tabularx}{0.8\textwidth}{>{\raggedright\arraybackslash}l >{\raggedright\arraybackslash}l >{\raggedright\arraybackslash}r >{\raggedright\arraybackslash}r}
%         \toprule
%         \textbf{Component} & \textbf{Usage/Quantity} & \textbf{A} & \textbf{B} \\
%         \midrule
%         Yada & yada & 1 & 2 \\
%         \midrule
%         \textbf{Total} & & \textbf{\(\sim\)\$0.00} & \textbf{\(\sim\)\$0.00} \\ 
%         \bottomrule
%     \end{tabularx}
% \end{table*}

% ======== References =========
\sloppy
\printbibliography[notcategory=selfref]
% \printbibliography

\end{document}